
\documentclass[a4paper,11pt]{article}

% --- Pakete ---
\usepackage[utf8]{inputenc}
\usepackage[T1]{fontenc}
\usepackage[ngerman]{babel}
\usepackage[left=2.5cm, right=2.5cm, top=2.5cm, bottom=2.5cm]{geometry}
\usepackage{xcolor}
\usepackage{enumitem}
\usepackage{titlesec}
\usepackage{parskip}
\usepackage{tcolorbox}
\usepackage{booktabs}
\usepackage{csquotes} % Für korrekte Anführungszeichen

% --- Design ---
\definecolor{rsBlue}{RGB}{0, 51, 102}
\definecolor{accentColor}{RGB}{204, 0, 0}
\definecolor{researchColor}{RGB}{70, 130, 180} % SteelBlue

\titleformat{\section}{\color{rsBlue}\Large\bfseries\uppercase}{}{0em}{}[\titlerule]
\titleformat{\subsection}{\color{rsBlue}\large\bfseries}{}{0em}{}

% Box für Forschungsfragen
\newtcolorbox{researchbox}[1]{
    colback=researchColor!10, 
    colframe=rsBlue, 
    title=\textbf{#1}, 
    fonttitle=\bfseries,
    boxrule=0.5mm,
    arc=2mm
}

\begin{document}

% --- KOPF ---
\begin{center}
    {\LARGE \textbf{Masterarbeit Proposal}}\\[0.8em]
    {\Large \textbf{Generative \enquote{Architecture is Code}: KI-basierte Synthese modularer SysML v2 Modelle aus heterogenen Bestandsdaten}}\\[2em]
\end{center}

\section{1. Motivation: Das \enquote{Schuhschachtel-Dilemma}}

Im Brownfield-Engineering stehen Unternehmen vor der Herausforderung, komplexe Bestandssysteme in modellbasierte Umgebungen (MBSE) zu überführen. Die Ausgangslage gleicht oft einer \enquote{Schuhschachtel}: Die Informationen liegen unstrukturiert, fragmentiert und in verschiedenen Formaten vor (CAD-Zeichnungen, Excel-Stücklisten, PDF-Spezifikationen).

Der aktuelle Ansatz der manuellen Modellierung führt zu zwei Problemen:
\begin{itemize}
    \item \textbf{Mangelnde Modularität (Silo-Problem):} Modelle werden oft \textit{ad-hoc} basierend auf den gerade verfügbaren Daten erstellt. Tauchen später neue Daten auf (z.\,B. eine Schnittstelle zu einem Nachbarsystem), passt das ursprüngliche Modell nicht mehr und muss aufwändig refaktoriert werden.
    \item \textbf{Verlust impliziten Wissens:} Ein erfahrener Ingenieur weiß, dass ein Motor immer eine Energieversorgung braucht, auch wenn dies nicht in der Stückliste steht. Einer \enquote{dummen} Übersetzung fehlt dieses Wissen, was zu unvollständigen Modellen führt.
\end{itemize}

\section{2. Forschungsleitende Fragestellungen}

Ziel der Arbeit ist es, dieses Dilemma durch einen KI-gestützten \enquote{Architecture is Code} (AiC) Ansatz zu lösen. Im Zentrum steht dabei nicht die einfache Übersetzung, sondern die intelligente Architektur-Synthese.

\begin{researchbox}{Hauptforschungsfrage (HF)}
    Wie lässt sich eine KI-gestützte Methodik gestalten, die aus heterogenen, unstrukturierten Bestandsdaten valide und modular erweiterbare SysML v2 Systemarchitekturen generiert?
\end{researchbox}

Zur Beantwortung der HF werden folgende Teilforschungsfragen (TF) adressiert:

\textbf{TF 1 (Antizipation \& Erwartungshaltung):}
Inwieweit können durch trainierte Architektur-Muster (Patterns) notwendige Strukturelemente (insb. Ports und Interfaces) antizipiert werden, die in den Quelldaten (noch) nicht explizit enthalten sind? 
\textit{(Vermeidung von \enquote{Dead-End}-Architekturen durch KI-Implikation).}

\textbf{TF 2 (Modularität \& Puzzle-Piece-Design):}
Welche Anforderungen muss der generierte SysML v2 Code erfüllen, um eine nachträgliche, non-destruktive Integration isoliert generierter Teilmodelle zu gewährleisten?
\textit{(Sicherstellung der \enquote{Plug \& Play}-Fähigkeit).}

\section{3. Methodischer Ansatz}

Das Lösungskonzept \enquote{AiC-Generator} basiert auf einem dreistufigen Prozess, der die genannten Forschungsfragen operativ umsetzt.

\subsection*{Phase A: Feature Ingestion (Die Schuhschachtel)}
Das System akzeptiert beliebige Datenfragmente als \enquote{Features}. 
\begin{itemize}
    \item Beispiel Input: \enquote{Bestellliste.xlsx} enthält \textit{Pumpe A} und \textit{Ventil B}.
    \item Herausforderung: Die Beziehung zwischen A und B ist im Excel oft nicht definiert.
\end{itemize}

\subsection*{Phase B: Generative Pattern Matching (Die Intelligenz)}
Hier greift TF 1. Anstatt starrer \enquote{Wenn-Dann}-Regeln nutzt die KI (LLM) generalisierte Erwartungshaltungen.
\begin{itemize}
    \item \textbf{Erwartung:} \enquote{Eine Pumpe in einem fluidischen System benötigt typischerweise einen Fluid-In, Fluid-Out und Power-In.}
    \item \textbf{Aktion:} Die KI generiert diese Ports als \enquote{Pending Interfaces}, auch wenn das Excel sie verschweigt. Dies bereitet das Modell auf zukünftige Verbindungen vor.
\end{itemize}

\subsection*{Phase C: Modular Code Synthesis (Der Output)}
Hier greift TF 2. Der generierte SysML v2 Code wird so strukturiert (Nutzung von Namespaces, Usages und Part Definitions), dass er wie ein Puzzleteil fungiert.
\begin{itemize}
    \item \textbf{Ergebnis:} Ein \enquote{Stand-Alone}-Modell, das syntaktisch korrekt ist, aber offene Schnittstellen für spätere Erweiterungen bereithält.
\end{itemize}

\section{4. Geplante Validierung (Proof of Concept)}

Die Methodik wird anhand eines \enquote{Blind-Tests} validiert:
\begin{enumerate}
    \item \textbf{Szenario 1:} Generierung eines Teilmodells A (z.\,B. \textit{Mechanik}) aus Datensatz 1.
    \item \textbf{Szenario 2:} Generierung eines Teilmodells B (z.\,B. \textit{Elektronik}) aus Datensatz 2 (ohne Kenntnis von A).
    \item \textbf{Integrationstest:} Prüfung, ob sich Modell A und Modell B in einem SysML v2 \texttt{root}-Modell verbinden lassen, ohne den Code von A oder B umschreiben zu müssen.
\end{enumerate}

\section{5. Erwarteter Erkenntnisgewinn}

Die Arbeit soll aufzeigen, dass der Einsatz von KI im Systems Engineering über die reine Assistenz (Chatbot) hinausgehen kann und als \textbf{Architektur-Enabler} fungiert. Sie liefert einen Blaupause für den Umgang mit Altdaten im Zeitalter von \enquote{Everything is Code}.

\end{document}